% IS 617 report: 4-5 pages + unlimited references and appendices, ACL template.
% The "% Rubric" comments quote the course grading rubric (points in brackets);
% see README.md in this folder. They are comments and do not appear in the PDF.
%
% The notes were rewritten on 2026-10-08, when the project moved from hiring to bank
% customer checks. The notes for the hiring design are in git history (commit 272792d).
\documentclass[11pt]{article}

\usepackage[preprint]{acl}

\usepackage{times}
\usepackage{latexsym}
\usepackage[T1]{fontenc}
\usepackage[utf8]{inputenc}
\usepackage{microtype}
\usepackage{inconsolata}
\usepackage{graphicx}
\usepackage{booktabs}
\usepackage{xcolor}

% --- Team notes ---------------------------------------------------------------
% The blue notes are visible in the PDF while drafting. They are facts and
% reminders, NOT report text: replace each one with the team's own wording.
% To hide all of them at once, change \shownotestrue to \shownotesfalse.
% They must be gone before submission.
\newif\ifshownotes
\shownotestrue
\newcommand{\teamnote}[2]{\ifshownotes\par\medskip\noindent{\small\color{blue}%
  \textbf{Notes: #1} (not report text; rewrite in your own words)\par #2\par}\medskip\fi}

\title{Which Arab? How LLMs Rate Bank Customers from the 22 Arab League Countries in Money-Laundering Checks}

\author{
  Laith Sandouk \and Aleksander Kasak \and Nico Fedotov \\
  University of Mannheim \\
  \texttt{\{laith.sandouk, aleksander.kasak, nico.fedotov\}@students.uni-mannheim.de}
}

\begin{document}
\maketitle

\teamnote{title}{%
The title above is a working title proposed by Claude Code on 2026-10-08 at Laith's
request. The team has to confirm or replace it. Other options, none of which assumes a
result:
\begin{itemize}
  \item Same File, Different Passport: Auditing LLMs in Bank Customer Checks
  \item One Label, 22 Nationalities: Do LLMs Tell Arab Bank Customers Apart?
  \item Which Passport Is High Risk? A Nationality-Swap Audit of LLMs in
        Anti-Money-Laundering Checks
\end{itemize}}

\begin{abstract}
% Rubric [5]: 250-word summary of your project including a formulation of its
% research questions and an overview of methods and results.
\teamnote{abstract}{%
At most 250 words. The rubric asks for the research questions, the methods and the
results. In all notes, (F) after a source means the full text or the cited passages were
read and (A) means only the abstract was read: check the full text before citing a
number from it. (press) is a press report and (vendor) a company's own statement.
\begin{itemize}
  \item Context: banks must check every new customer for money-laundering risk. AI
        systems are being sold for this check, and EU law regulates their use in it.
  \item Why Arab customers: LLMs repeat stereotypes about Arabs and Muslims in what they
        write. Banks in Europe answered enquiries from Arabic-sounding names far less
        often. Numbers and sources: introduction notes.
  \item Why this task: the nationality cannot be left out of the file (the bank must
        record it), and the rules say how far it may count (they speak of residence, not
        of the passport).
  \item Gap: every study treats ``Arab'' as one group. We found no study in which an LLM
        rates bank customers while their nationality is varied; this rests on a few web
        searches only.
  \item Questions (from the proposal of 2026-10-08; hypotheses are not written yet):
        (1) with the customer file held fixed, do a model's ratings differ across the 22
        Arab League nationalities, and by more than across unrelated nationalities?
        (2) Do the differences follow the EU list of high-risk countries? (3) Arab League
        nationalities against German, Turkish and placebo nationalities. (4) Does a
        paragraph saying that nationality is not a risk factor change this? (5) Do the
        four models agree?
  \item Method: 20 synthetic customer files; 28 nationalities, changed in the nationality
        and passport lines only; a one-page bank policy in two versions; risk score
        0--100, rating, enhanced checks, recommendation; four open-weight models.
  \item Results: nothing has been run as of 2026-10-08. Write this part last, and only
        from the blocks that were run.
  \item Last sentence: what follows. The data decide: the models separate the
        nationalities, lump them together, or follow the list.
\end{itemize}}
\end{abstract}

\section{Introduction}
\label{sec:intro}
% Rubric "Motivation" [5]: What question(s) do you seek to answer and why? Elaborate on
% how this research matters in terms of science, technology and/or society.

\teamnote{motivation, step by step}{%
\begin{itemize}
  \item \textbf{1. LLMs in decisions about people.} Audits mostly vary race and gender,
        often through names. Nationality is rare: it appears in 13.2\% of 189 papers on
        LLM bias, gender in 79.9\%. In finance: one line stating the applicant's race on
        a mortgage file; GPT-4 Turbo recommended approval 8.5 percentage points less
        often and an interest rate 35 basis points higher for Black applicants. Sources:
        \texttt{ghosh2025bias} (A), \texttt{bowen2025measuring} (F, sections 1--4).
  \item \textbf{2. LLMs and Arabs.} In generated text: Muslims tied to violence in 66 of
        100 GPT-3 completions; MENA names 78\% under-represented in generated workplace
        stories; Middle Eastern personas described through religion. Models also tell
        Arab countries apart: Libyan, Sudanese and Tunisian among the five most negative
        demonyms in GPT-2 text; additional scrutiny suggested for applicants from Iraq,
        Sudan and Somalia in simulated immigration decisions, while most country effects
        in the choices were not significant; the values of some Arab countries
        represented much better than others. Sources: \texttt{abid2021persistent} (F),
        \texttt{shieh2026intersectional} (F), \texttt{cheng2023marked} (F),
        \texttt{venkit2023nationality} (F), \texttt{mao2025gatekeepers} (F),
        \texttt{zahraei2025menavalues} (F).
  \item \textbf{3. Banks.} Field experiment, 1,218 e-mails to banks in seven European
        countries, asking only for contact details: 55.2\% replies to domestic-sounding
        names, 31.6\% to Arabic-sounding names; loan enquiries 59.8\% vs 35.1\%. The
        Netherlands Institute for Human Rights ruled in 2024 that ING discriminated
        against two customers whose payments were blocked and checked because of names
        that did not sound Dutch. Campaigners say HSBC closed dozens of accounts of
        Syrian nationals living in the UK from 2014. Sources:
        \texttt{stefan2018ethnical} (F; the list of countries is in an appendix that was
        not read, so do not write that Germany is among them),
        \texttt{nltimes2024ing} (press), \texttt{mee2016hsbc} (press).
  \item \textbf{4. One category.} All of these treat ``Arab'' or ``Arabic-sounding'' as
        one group. A name cannot do better: raters assigned 34\% of Moroccan names to
        Morocco. Source: \texttt{martiniello2022signaling} (F).
  \item \textbf{5. Why a bank's customer check.} A German bank must record name, place
        and date of birth, nationality and address of every customer, and rate the
        money-laundering risk. So the nationality is always in the file. Sources:
        \texttt{gwg2017} (F, sections 11 and 15, annex 2).
  \item \textbf{6. The yardstick.} European guidelines: country risk is about where the
        customer is based or resident, the main places of business, and relevant links.
        The word ``nationality'' does not occur in the 225 pages of the final report. An
        isolated risk factor does not necessarily move a customer into a higher category.
        The Commission's list of high-risk third countries has 26 entries; four are Arab
        League members: Algeria, Lebanon (since 5 August 2025), Syria, Yemen (since
        23 September 2016). German law applies the list to relationships and transactions
        involving such a country or a person resident there. No disadvantage on account
        of nationality when opening a payment account. No blanket refusal of categories
        of customers. Sources: \texttt{eba2021riskfactors} (F, guidelines 2.9 and 3.3,
        word search), \texttt{ec2026highrisk} (F, read 8 October 2026),
        \texttt{gwg2017} (F), \texttt{zkg2016} (F, section 3),
        \texttt{eba2023derisking} (F, paragraphs 9--10).
  \item \textbf{7. AI in this task.} EU regulation 2024/1624, Article 76(5), applies from
        10 July 2027: decisions ``resulting from automated processes [\ldots] or from
        processes involving AI systems'' are allowed if accepting or refusing a customer
        is ``subject to meaningful human intervention'' and the customer can get an
        explanation. Vendors sell AI agents for customer checks and screening alerts;
        one says 10 of the top 20 banks use its agents. ECB: strong increase in AI use
        cases in banks between 2023 and 2024; for credit scoring and fraud detection
        mainly decision trees. Sources: \texttt{eu2024amlr} (F, Articles 76(5) and 90),
        \texttt{workfusion2026agents} (vendor), \texttt{sardine2026agentic} (vendor),
        \texttt{ecb2025ai} (F).
  \item \textbf{8. Gap and what we do.} No study was found in which an LLM rates bank
        customers while their nationality is varied. This rests on a few web searches of
        8 October 2026, not on a systematic search: write ``to our knowledge'' only after
        one has been done. Then the question and one sentence on the design.
\end{itemize}}

\teamnote{research questions (as in the proposal of 2026-10-08)}{%
\begin{itemize}
  \item Framing: the 22 countries differ in wealth, stability and how their citizens are
        seen abroad. Does a model carry that into a decision about an individual
        customer?
  \item RQ1 (main): do risk ratings and recommendations differ across the 22 Arab League
        nationalities with the file otherwise identical, and more than across the
        placebo nationalities?
  \item RQ2: do the differences follow the EU list of high-risk countries (four of the
        22), or a country's economic and political position?
  \item RQ3: are Arab League nationalities as a group rated differently from German,
        Turkish and placebo nationalities?
  \item RQ4: do the differences remain when the policy says that nationality is not a
        risk factor?
  \item RQ5: do the four models agree?
\end{itemize}}

\teamnote{why it matters (rubric: science, technology, society)}{%
\begin{itemize}
  \item Science: bias audits report one number for ``Arab''. Here there is an official
        yardstick for the 22 countries, so three outcomes can be told apart: the model
        follows the list, lumps all 22 together, or draws its own lines.
  \item Technology: AI agents are sold for customer checks (step 7), and EU law expects
        AI systems in this decision. How language models react to the nationality in a
        customer file is untested, as far as we found.
  \item Society: an account is needed for wages, rent and daily life. A higher risk
        rating means more checks, delays or a refusal. This last sentence is our own
        argument and has no source.
\end{itemize}}

\teamnote{do not overstate}{%
\begin{itemize}
  \item Not ``banks use LLMs to decide on customers'': what the products do, which
        models they use and which fields the model sees is not public. We test open
        models in the task the products are sold for.
  \item Not ``LLMs discriminate against Arab customers'' as a fact. Nothing has been run.
        Explicit origin labels can also make models more cautious
        (\texttt{bunel2026customer}, F).
  \item A higher rating for a listed country is not automatically wrong in practice.
        Our argument is narrower: for a customer who lives and works in Germany and has
        no payments abroad, the written rules give the passport no weight.
  \item The legal texts were read in single sections only (see the literature matrix).
        Nobody on the team is a lawyer. Do not present the policy as legal advice.
  \item Say ``Arab League nationalities'', not ``MENA'' and not ``Arab customers''.
  \item The claim to be first needs a systematic search that has not been done.
  \item Do not cite numbers from \texttt{mohammad2026mirage}: they are placeholders.
  \item This is an audit of model behaviour, not a tool for checking customers. The risk
        score is not a probability.
\end{itemize}}

\section{Related Work}
\label{sec:related}
% Rubric "Research background" [5]: What previous work is related to this project and
% what can we learn from it? Do not just cite papers but highlight what's relevant to
% your work. Reflect about the cited contributions on your text.

\teamnote{what each strand teaches us}{%
The rubric asks for what is relevant to our work, not for a list of papers.
\begin{itemize}
  \item \textbf{LLM audits with one explicit line.} Mortgage files with a race line:
        gaps in approval and interest rate, larger for weak files; the rate gap in all
        eight models. Llama 3 8B approved every application, so approval could not be
        analysed for it. Claude 3 Opus answered only 74\% of the requests for Black
        applicants (\texttt{bowen2025measuring}, F). For us: a yes/no or three-level
        answer may not move in small models, hence the 0--100 score; count refusals by
        nationality.
  \item \textbf{Direction is unstable.} Real origin labels led to fewer discriminatory
        recommendations than invented groups (\texttt{bunel2026customer}, F); every
        detected race or ethnicity bias favoured the minority in one study
        (\texttt{arcuschin2026blind}, F). For us: expect either direction.
  \item \textbf{Explicit cues.} A stated attribute moves decisions more than a name
        (\texttt{tamkin2023evaluating}, F). Different cues for the same group can give
        different conclusions (\texttt{tonneau2026cues}, A). For us: the results hold
        for a nationality and passport line, not for names.
  \item \textbf{Nationality in LLM decisions.} 20 nationalities in job recommendations,
        Jordan the only Arab one (\texttt{salinas2023unequal}, F). GPT-4 prefers
        developers from some countries, tested by swapping only the location string
        (\texttt{nakano2024nigerian}, F). Country effects mostly not significant in
        simulated immigration choices (\texttt{mao2025gatekeepers}, F).
  \item \textbf{Arab group in LLM studies.} Always pooled:
        \texttt{lippens2024computer} (F), \texttt{hoffmann2026evaluating} (F),
        \texttt{bai2025explicitly} (A), \texttt{albaroudi2026addressing} (F). These are
        hiring studies; none is about banking.
  \item \textbf{Telling the model to ignore the attribute.} Brought discrimination close
        to zero in Claude 2.0 (\texttt{tamkin2023evaluating}, F). A ``use no bias''
        sentence removed the approval gap in mortgage files
        (\texttt{bowen2025measuring}, F). Instructions fail once realistic context is
        added in hiring (\texttt{karvonen2025robustly}, F). For us: the second policy
        version can go either way.
  \item \textbf{Models may notice a test.} \texttt{vohra2026audit} (F),
        \texttt{needham2025large} (A).
  \item \textbf{Human evidence in banking.} See the introduction notes, step 3.
  \item \textbf{Not searched yet.} Research on bias in anti-money-laundering systems and
        on ``de-risking''. Find sources, or say that this literature was not covered.
\end{itemize}}

\section{Data}
\label{sec:data}
% Rubric [10]: Describe your data sources including descriptive statistics and plots.
% Carefully document all references to data sources.

\teamnote{materials}{%
\begin{itemize}
  \item 20 customer files for a current account, one per occupation (the 20 occupations
        of the earlier hiring design). Net incomes are invented.
  \item Held constant in every file: lives in Germany, tax resident in Germany, employed
        by a German employer, salary as source of funds, no payments to or from other
        countries, no sanctions match, not politically exposed, no adverse media.
  \item 10 plain files, 10 with one mild indicator named in the policy: account opened
        by video call (4), cash deposits from tips (2), second income in cash (2),
        second income without cash (2).
  \item Nationality: 28 conditions. 22 Arab League nationalities, 2 benchmarks (German,
        Turkish), 4 placebo nationalities (Uruguayan, Malawian, Maldivian, Cambodian).
        No version without a nationality: the field is mandatory.
  \item What changes with the nationality: the nationality line and the passport line.
        All non-German customers hold an unlimited settlement permit. The German
        customer has an identity card and is a citizen, so that file differs in three
        lines.
  \item The policy: one page, a paraphrase of public rules, with the Commission's list
        of 26 high-risk countries as of 29 January 2026. Second version: one more
        paragraph saying that nationality is not a risk factor.
  \item Country data: EU list status in \texttt{customer\_nationalities.csv}; World Bank
        income and region in \texttt{country\_covariates.csv}.
  \item Everything is synthetic: no real customer, bank or employer. The bank's name
        (Lindenufer Bank) is invented.
  \item Still needed for the rubric: descriptive statistics and plots of the answers.
\end{itemize}
Source: the README files under \texttt{experiment/}.}

\teamnote{what a customer file looks like}{%
The file shows a reference number instead of a name. Example of the lines that change:
\begin{quote}\ttfamily
Nationality: Jordanian\\
Identity document: Jordanian passport, checked in the original\\
Residence status: Settlement permit (Niederlassungserlaubnis), unlimited
\end{quote}
\begin{itemize}
  \item A name cannot tell Arab countries apart, and it adds signals we do not want
        (gender, religion).
  \item A real file also holds the place of birth. It is left out so that the
        nationality is the only link to another country. Say so as a limitation.
\end{itemize}}

\teamnote{nationalities whose Arab identity is contested}{%
Comoros, Djibouti and Somalia are Arab League members whose Arab identity is contested.
They stay in every comparison. Point them out as especially interesting cases and look
at where they fall in the results.}

\section{Methods}
\label{sec:methods}
% Rubric [10]: Explain how you filtered data, normalized values, computed additional
% variables, etc. Detail the statistical analyses and other methods you will apply to
% assess your research question(s). Carefully document all references to methods and
% packages used.

\teamnote{experiment setup}{%
\begin{itemize}
  \item One run: a main prompt (the model assists the onboarding team of a bank in
        Germany; an analyst takes the final decision; the policy), then a message with
        one customer file and the request.
  \item Answer: a JSON object with a risk score 0--100, a rating (low, medium, high),
        enhanced checks yes/no, a recommendation (open, escalate, decline) and a short
        reason.
  \item Two policy versions: without and with the paragraph on nationality.
  \item Three wordings of the request.
  \item Five repetitions per prompt are planned; the decoding settings are not decided.
  \item Models: Qwen3 8B, Llama 3.1 8B, Gemma 3 12B, Mistral Small 3.2 24B. State the
        exact versions and settings.
  \item \texttt{scripts/build\_runs.py} builds and checks the prompts. It calls no
        model. The code that sends the runs does not exist yet.
  \item How the answers are analysed: not decided as of 2026-10-08.
\end{itemize}
Planned runs per model:\par\smallskip
\begin{tabular}{@{}lr@{}}
\toprule
Block & Runs \\
\midrule
1 Policy without the nationality paragraph & 8,400 \\
2 Policy with the nationality paragraph & 8,400 \\
\midrule
Total & 16,800 \\
\bottomrule
\end{tabular}\par\smallskip
Each block: 20 files $\times$ 28 nationalities $\times$ 3 wordings $\times$ 5
repetitions. Source: \texttt{experiment/README.md}.}

\section{Results}
\label{sec:results}
% Rubric [10]: Present the results of your analysis including tables and figures that
% communicate and illustrate those results.

\teamnote{nothing to report yet}{%
As of 2026-10-08 no run exists. The proposal lists what to compare:
\begin{itemize}
  \item the spread of risk scores and ratings across the 22 Arab League nationalities,
        against the spread across the 4 placebo nationalities;
  \item the four listed countries against the eighteen others;
  \item unlisted Arab League nationalities against placebo, German and Turkish;
  \item with versus without the paragraph on nationality;
  \item plain files against files with one indicator;
  \item whether the four models agree;
  \item how often a model refuses or gives an unreadable answer, by nationality.
\end{itemize}
The hypotheses and the analysis have to be written down before the first real run.
A pilot has to show first whether the answers vary at all. Block 1 on all four models is
the plan for the midterm.}

\section{Conclusion and Discussion}
\label{sec:conclusion}
% Rubric [5]: Evaluate answers to the question and their reliability. Identify
% limitations and alternative explanations for your results.

\teamnote{alternative explanations to discuss}{%
\begin{itemize}
  \item A passport also says ``foreigner''. The German file is the only one of a
        citizen, and it differs in three lines. Compare the Arab League nationalities
        with the Turkish and placebo files, which are foreigners too.
  \item The policy names the listed countries. A model may apply a rule about residence
        to a passport because the country name appears in the prompt. A run without the
        list would separate this; it is not planned.
  \item The label may stand for religion, conflict or income of the country, not for
        the nationality as such. \texttt{country\_covariates.csv} has income and region.
  \item Any set of labels spreads the scores a little. The placebo nationalities are the
        yardstick; with only four of them it is a rough one.
  \item Models may notice that this is a test (\texttt{vohra2026audit}, F;
        \texttt{needham2025large}, A).
  \item All 20 customers are ordinary. If a model answers ``low'' for everybody, the
        rating cannot show a difference; the score may.
  \item A result of ``no difference'' is an answer too.
\end{itemize}}

\section*{Limitations}

\teamnote{limits to state}{%
\begin{itemize}
  \item FILL IN AFTER THE RUNS: which blocks and which models were actually run. As of
        2026-10-08 nothing has been run.
  \item The products in this field run on large commercial models and their prompts are
        not public. We test open models of 8 to 24 billion parameters with a policy and
        a file that we wrote.
  \item In practice most customer risk ratings come from fixed scorecards; language
        models review cases and write the reasoning. We imitate that step.
  \item The policy is a paraphrase by non-lawyers of rules read in part.
  \item English prompts for a German bank; one bank, one product, one country.
  \item Synthetic files without name and place of birth; nationality given as explicit
        lines.
  \item One list (the Commission's) on one date. Other lists differ and change.
\end{itemize}}

\section*{Author Contributions}
% Rubric: after the conclusion; does not count towards the page limit; as detailed as
% possible up to a maximum of 1 page; follow the CRediT guidelines.

\section*{Use of AI Tools}
% Rubric: note which tasks have been assisted by AI tools (including identifying the
% tools) at the end of the report.

\teamnote{what to declare}{%
As of 2026-10-08, assisted by Claude Code: literature search and research notes; the
search for and comparison of scenarios (\texttt{research/scenario\_reviews/}); design
documents; the bank policy text, the customer files and the prompts; the first drafts of
both project proposals (2026-10-04 and 2026-10-08); repository code and structure; the
working title and title suggestions; and these notes. For the earlier hiring design also
the job ads, CVs and prompts, which are no longer used. Keep the list up to date.}

\section*{Declaration}
% AI usage policy: a declaration at the end stating that you have written the whole
% report on your own, i.e., without any external help.

% Uncomment once the text cites something; an empty bibliography fails to compile.
% \bibliography{custom}

\appendix

\section{Prompts}
\label{sec:appendix-prompts}

\teamnote{prompts}{%
Copy in, unchanged, the main prompt, both versions of the policy and one wording of the
request from \texttt{experiment/prompts/} and \texttt{experiment/policy/}, with one
customer file.}

\end{document}
